
\subsection*{1.3.}
Rappelons (IV 2.10) que le bifoncteur
\[
(\underline{E},\underline{F})\rightsquigarrow \underline{E}\otimes_{\underline{O}_X}\underline{F},
\]
où $\underline{E}$ est un $\underline{O}_X$-module localement libre et $\underline{F}$ un $\underline{O}_X$-module cohérent, induit une loi de composition
\begin{equation}\tag{1.3.1}
K^{\bullet}(X)\times K_\bullet(X)\longrightarrow K_\bullet(X)
\end{equation}
qui fait de $K_\bullet(X)$ un $K^{\bullet}(X)$-module unitaire. Le résultat-clé du présent exposé est le

\textbf{Théorème 1.3.2.} \uline{Soit} $X$ \uline{un schéma noethérien. Pour tout} $i,j\in\mathbb Z$ \uline{la loi de composition} (1.3.1) \uline{induit par restriction une loi de composition}
\begin{equation}\tag{1.3.3}
\Filt^i(X)\times\Filt_j(X)\longrightarrow\Filt_{j-i}(X).
\end{equation}

\textbf{Corollaire 1.3.3.} \uline{Si} $X$ \uline{est un schéma noethérien de dimension} $\leq j$\uline{, alors l'homomorphisme canonique}
\[
\Theta_X:K^{\bullet}(X)\longrightarrow K_\bullet(X)
\]
\uline{défini par} $\Theta_X(x)=x\cdot\cl_\bullet(\underline{O}_X)$ ($x\in K^{\bullet}(X)$\uline{) induit pour tout} $i\in\mathbb Z$ \uline{un homomorphisme}
\[
\Filt^i(X)\longrightarrow\Filt_{j-i}(X).
\]

\emph{Démonstration :} Montrons d'abord que (1.3.3), pour $i+j\leq k\in\mathbb Z$ et pour \emph{tout} sous-schéma intègre fermé de dimension $\leq j$, implique (1.3.2) pour $i+j\leq k$. Grâce à (1.1.4) il suffit de montrer que pour tout $x\in\Filt^i(X)$ et pour \emph{tout} sous-schéma intègre fermé $Z$ de dimension $\leq j$ on a
